\section{Application of Program Outcomes}
Table~\ref{tab:po} maps each Program Outcome and Program Specific Outcome addressed by this case study to its concrete implementation. Every mapping demonstrates how foundational engineering knowledge, modern tools, and professional ethics are exercised within the project.

{\small
\begin{longtable}{@{}>{\centering\arraybackslash}p{1.8cm} p{13.2cm}@{}}
    \caption{Mapping of Program Outcomes and Program Specific Outcomes to the Case Study.}\label{tab:po}\\
    \toprule
    \textbf{Outcome} & \textbf{Outcome Description and Concrete Relevance to the Case Study} \\
    \midrule
    \endfirsthead
    \toprule
    \textbf{Outcome} & \textbf{Outcome Description and Concrete Relevance to the Case Study} \\
    \midrule
    \endhead
\textbf{PO1} & \textbf{Engineering Knowledge:} Applies principles of mathematics, computer science, and engineering fundamentals to formalize, model, and analyze the system architecture, mathematical constraints, and foundational algorithms. In this case study (Industrial Equipment Monitoring Using IoT), transduction physics, link budgets, and embedded digital signal processing are mathematically derived for Industrial Equipment Monitoring Using IoT. \\
    \addlinespace
\textbf{PO2} & \textbf{Problem Analysis:} Systematically identifies, formulates, and investigates complex computational challenges, diagnosing architectural bottlenecks, concurrency anomalies, latency constraints, and integrity failures through rigorous analytical techniques. In this case study (Industrial Equipment Monitoring Using IoT), systematic diagnosis of RF attenuation, battery exhaustion, and sensor drift establishes resilient engineering constraints. \\
    \addlinespace
\textbf{PO3} & \textbf{Design/Development of Solutions:} Designs scalable system architectures, robust database schemas, object hierarchies, or automated pipelines meeting specified operational needs with proper consideration for efficiency, maintainability, and safety. In this case study (Industrial Equipment Monitoring Using IoT), end-to-end hardware-software architecture with ultra-low-power sleep states, fail-safe relays, and cloud brokers is developed. \\
    \addlinespace
\textbf{PO5} & \textbf{Modern Tool Usage:} Selects and applies state-of-the-art software frameworks, modeling tools, database management engines, profiling utilities, and development environments to deliver high-performance solutions. In this case study (Industrial Equipment Monitoring Using IoT), modern embedded IDEs, MQTT protocol analyzers, InfluxDB time-series engines, and Grafana dashboards are utilized. \\
    \addlinespace
\textbf{PO7} & \textbf{Environment and Sustainability:} Analyzes the environmental impact, compute resource consumption, and sustainability profiles of the deployed computational processes and operational paradigms. In this case study (Industrial Equipment Monitoring Using IoT), ultra-low-power hardware design and ecological resource optimization minimize carbon footprints and electronic waste. \\
    \addlinespace
\textbf{PSO1} & \textbf{Program-Specific (AI/ML Core):} Leverages advanced intelligent computational techniques, algorithmic optimization, predictive analytics, and automated decision-making workflows to enhance system capabilities. In this case study (Industrial Equipment Monitoring Using IoT), intelligent data collection, feature extraction, and edge analytics feed directly into machine learning pipelines. \\
    \addlinespace
\textbf{PSO2} & \textbf{Program-Specific (Applied Engineering):} Designs and implements production-grade end-to-end software pipelines integrating intelligent modules with robust backend infrastructure and enterprise services. In this case study (Industrial Equipment Monitoring Using IoT), seamless integration of embedded firmware, industrial sensors, and enterprise cloud dashboards fulfills applied engineering standards. \\
    \addlinespace
    \bottomrule
\end{longtable}}